\chapter{Homework 9 (due 11/13)}

\subsection{Coherent states}

Consider a normalized coherent state $\dket{z}$ of a simple harmonic oscillator of mass $m$ and frequency $\omega$, defined as follows:
\begin{align}
\dket{z}\equiv e^{-|z|^2/2}e^{za^\dagger}\dket{0}=e^{-|z|^2/2}\sum_{n=0}^\infty\frac{z^n}{\sqrt{n!}}\dket{n},~~~z\in\mathbb{C}.
\end{align}

(1) Show that the coherent state defined above is an eigenstate of annihilation operator $\hat a$:
\begin{align}
\hat a\dket{z}=z\dket{z}
\end{align}

(2) Show that two different coherent states are not orgononal, with the following inner product:
\begin{align}
\expval{z|z^\prime}=e^{z^\ast z^\prime-|z|^2/2-|z^\prime|^2/2},~~~z,z^\prime\in\mathbb{C}.
\end{align}

(3) \textbf{ Bonus} Prove the completeness relation for coherent states
\begin{align}
\hat 1=\int\frac{\text{d}^2z}{\pi}\dket{z}\dbra{z}
\end{align}
Hint: for a complex variable $z=x+\imth y\in\mathbb{C}$, we have $\text{d}^2z=\text{d}x~\text{d}y$.

Comment: since coherent states are not orthogonal to each other, they form an overcomplete set of basis for the Hilbert space. 

(4) Show that under the time evolution by a simple hamonic oscillator Hamiltonian, a coherent state evolves as
\begin{align}
e^{-\imth\hat H_\text{SHO}t/\hbar}\dket{z}=e^{\imth\phi(t)}\dket{ze^{-\imth\omega t}}
\end{align}
Find out the phase $\phi(t)$. 

(5) A displacement operator $D(z)$ is defined as
\begin{align}\label{displacement operator}
D(z)=e^{za^\dagger-z^\ast a},~~z\in\mathbb{C}
\end{align}
Show that $D(z)$ is a unitary operator, and further show that
\begin{align}
D(z)\dket{0}=\dket{z}
\end{align}
where $\dket{0}$ is the normalized ground state of a SHO. 

Hint: according to the Baker-Campbell-Hausdorff formula, if 
\begin{align}
[\hat A,\hat B]=\alpha\hat 1
\end{align}
we have
\begin{align}
e^{\hat A+\hat B+\frac12[\hat A,\hat B]}=e^{\hat A}e^{\hat B}
\end{align}

\subsection{Quantization of electromagnetic fields in a one-dimensional cavity}

(1) Considering a propagating mode in a one-dimensional optical cavity, with momentum $k$ and frequency $\omega=c|k|$, its electric field is given by
\begin{equation}
 \hat E_x(y,t)=
 \sqrt{\frac{\hbar\omega}{\epsilon L}}
 \left(\hat a e^{\imth(ky-\omega t)}
 -\hat a^\dagger e^{\imth(\omega t-ky)}\right).
\end{equation}
For phase factor defined below
\begin{equation}
\Phi=\omega t-ky,
\end{equation}
show the electric field can be written as 
\begin{equation}
\hat E_x(y,t)=\sqrt{\frac{2\hbar\omega}{\epsilon L}}
\left(\hat X\cos\Phi+\hat Y\sin\Phi\right).
\end{equation}
where the quadratures $X,Y$ are defined as
\begin{equation}
\hat X=\frac{\hat a+\hat a^\dagger}{\sqrt2},
\qquad
\hat Y=\frac{\imth(\hat a^\dagger-\hat a)}{\sqrt2}.
\end{equation}

(2) In classical electromagnetism, the energy density is given by
\begin{align}
    u=\frac{\epsilon}{2}|\textbf{E}|^2+\frac{1}{2\mu}|\textbf{B}|^2
\end{align}
where $\epsilon$ is the dielectric constant and $\mu$ is the permeability. For a single mode of standing wave in the one-dimensional optical cavity (along $y$ direction), as we discussed in section \ref{section:1d optical cavity}, the electric and magnetic fields are given by
\begin{equation}
 {
 \hat E_x(y,t)=\imth
 \sqrt{\frac{\hbar\omega}{\epsilon L}}
 \sin(ky)
 \left(\hat a e^{-\imth\omega t}
 -\hat a^\dagger e^{\imth\omega t}\right).}
\end{equation}
and
\begin{equation}
{
 \hat B_z(y,t)=-
 k\sqrt{\frac{\hbar}{\epsilon L\omega}}
 \cos(ky)
 \left(\hat a e^{-\imth\omega t}
 +\hat a^\dagger e^{\imth\omega t}\right).}
\end{equation}
Show that the total energy of electromagnetic fields is carried by the photon:
\begin{align}
    \int_0^L u(y,t)\dif{y}=\hbar\omega (\hat a^\dagger\hat a+\frac12)
\end{align}


(3) In classical eletromagnetism, the momentum density of the electromagnetic fields is given by 
\begin{align}
    \textbf{g}=\epsilon~\textbf{E}\times\textbf{B}
\end{align}
Again we consider a single mode of standing wave in the one-dimensional cavity with $g_y=-\epsilon E_xB_z$, and compute total momentum $\int_0^L g_y(y,t)$ of electromagnetic fields carried by this standing wave. 


(4) In this problem we discuss the Planck formula of blackbody radiation, by considering electromagnetic fields as standing waves in a 3-dimensional cube of size $L\times L\times L$. The total energy density of photons (including the two polarization modes) is given by
\begin{align}
    u=\frac2{2\pi^2}\int|k|^2\dif|k|\expval{\hat n_k}\omega_k,~~~\omega_k=c|k|.
\end{align}
where the expectation value of photon number $\hat n_k=\hat a_k^\dagger\hat a_k$ is for thermal equilibrium at temperature $T>0$. By changing the integration variable from $|k|$ to (angular) frequency $\omega=c|k|$, derive the Planck formula:
\begin{align}
    u=\int\dif\omega~u(\omega),~~~u(\omega)=\frac1{\pi^2c^3}\frac{\omega^3}{e^{\beta\hbar\omega}-1}
\end{align}

Comment: in the original version of Planck's formula discussed in section \ref{sec:blackbody radiation}, we used the light frequency $\nu=\omega/2\pi$, hence the extra factor of $(2\pi)^4$. 


\subsection{Casimir effect}

Consider the quantization of light in a one-dimensional optical cavity of length $L$, as we discussed in section \ref{section:1d optical cavity}. 

(1) Compute the ground state energy $E_0$ (i.e. zero point energy) of the system as a function of $L,c,\hbar$. 

Hint: by the zeta function regularization we have
\begin{align}
    \sum_{n=1}^\infty n=\zeta{(-1)}=-\frac1{12}
\end{align}

(2) Show that $E_0$ increases with an increasing cavity length $L$, and compute the attractive Casimir force between the two cavity walls (separated by a distance $L$): 
\begin{align}
    F_\text{Cas}=-\frac{\partial E_0}{\partial L}
\end{align}




\subsection{(Bonus) Squeezed states}

We consider a generic squeezed state $\dket{\xi,z}$, characterized by two complex numbers $\xi,z\in\mathbb{C}$ as
\begin{align}\label{squeezed state}
\dket{\xi,z}\equiv\hat S(\xi)\hat D(z)\dket{0}\propto\hat S(\xi)\dket{z}
\end{align}
where $\dket{0}$ is the normalized ground state of a dimensionless (set $\hbar=m=\omega=1$) simple harmonic oscillator
\begin{align}
\hat H=\frac{\hat x^2+\hat p^2}{2}=a^\dagger a+\frac12
\end{align}
with 
\begin{align}\label{SHO:dimensionless}
[\hat x,\hat p]=\imth,~~~a=\frac{\hat x+\imth\hat p}{\sqrt2}e^{\imth\frac\theta2},~a^\dagger=\frac{\hat x-\imth\hat p}{\sqrt2}e^{-\imth\frac\theta2},~~~\forall~\theta\in\mathbb{R}.
\end{align}
The unitary squeezing operator $\hat S(\xi)$ is defined as
\begin{align}\label{suqeezing operator}
\hat S(\xi)\equiv e^{\frac{\xi^\ast a^2-\xi(a^\dagger)^2}2},~~~\xi\in\mathbb{C}
\end{align}
while the displacement operator $\hat D(z)$ is defined in (\ref{displacement operator}).

(1) \textbf{ Bonus} Show that for $\xi=re^{\imth\theta}$ with $r,\theta\in\mathbb{R}$, we have
\begin{align}
&S^\dagger(\xi)aS(\xi)=a\cosh{r}-e^{\imth\theta}a^\dagger\sinh r,\\
&S^\dagger(\xi)a^\dagger S(\xi)=a^\dagger\cosh{r}-e^{-\imth\theta}a\sinh r
\end{align}

Hint: use the Baker-Campbell-Hausdorff formula
\begin{align}
e^{A}Be^{-A}=B+[A,B]+\frac{1}{2!}[A,[A,B]]+\frac{1}{3!}[A,[A,[A,B]]]+\cdots
\end{align}

(2) \textbf{Bonus} Show that the squeezed state (\ref{squeezed state}) satisfies:
\begin{align}\label{squeezed state:gaussian wave packet}
(\hat x+\imth\hat pe^{-2r}-C)\dket{\xi,z}=0.
\end{align}
where position and momentum operators $\hat x,\hat p$ are defined in (\ref{SHO:dimensionless}), and $C\in\mathbb{C}$ is a constant. 

Hint: $(a-z)\dket{z}=0$ for a coherent state $\dket{z}$, which is related to the squeezed state by the squeezing operator $S(\xi)$ as shown in (\ref{squeezed state}). 

Comment: Eq. (\ref{squeezed state:gaussian wave packet}) shows that each squeezed state $\dket{\xi,z}$ is a Gaussian wave packet, with a characteristic radius of $R=e^{-r}$. 

(3) \textbf{ Bonus} Now let's consider the electric field operator in single-mode quantum optics
\begin{align}
\hat E(\chi)=\frac{a e^{-\imth\chi}+a^\dagger e^{\imth\chi}}{\sqrt2}
\end{align}
where 
\begin{align}
\chi=\omega t-{\vec k}\cdot\vec r+\pi/2
\end{align}
is the phase factor. Compute its variance 
\begin{align}
\Delta E^2(\chi)\equiv\expval{\xi,z|\big[\Delta\hat E(\chi)\big]^2|\xi,z}
\end{align}
w.r.t. the squeezed state (\ref{squeezed state}). For the specific case of $\xi=2$, plot $\Delta E^2(\chi)$ as a function of $\chi$. 
